% Independent mathematical proof section for the PA194H5 release.
% Requires amsmath, amssymb, amsthm; theorem and lemma environments.
\section{A PA proof of the consistency calibration for PA194H5}
\label{sec:h5-calibration}

The construction combines a coefficient-five formula encoding with the
visible-result, no-push MP operation and a Shift operation that tests only
the cached formula. The logical calibration below is an ordinary
mathematical proof of a PA-provability statement. Its finite machine facts
are the delivered source, transition, symbolic-runtime and blank-start
certificates. It is not a newly compiled end-to-end Lean or Rocq theorem.

Let $M$ denote the exact delivered dense table, with start zero and halt
label 194, and let $H_M$ assert that it has a finite halting trace from the
blank binary tape. The statement to be proved is
\[
 \mathrm{PA}\vdash
 H_M\ \longleftrightarrow\
 \exists p\,\operatorname{Prf}_{\mathrm{PA}}(p,\ulcorner0=1\urcorner).
 \tag{H5}
\]
Negating this equivalence in PA gives the requested consistency statement.
No consistency assumption is used to prove (H5).

\subsection{The total numerical calculus}

Put $P(a,b)=(a+b)(a+b+1)/2+b$ and define
\[
\begin{aligned}
 I_5(A,B)&=5P(A,B),&M_5(a,b)&=5P(a,b)+1,\\
 U_5(a,B)&=5P(B,a)+2,&E_5(a,b)&=5P(a,b)+3.
\end{aligned}
\]
The order of the two components in $U_5$ is intentional. Zero denotes
truth and $1=M_5(0,0)$ denotes the target false formula. Write
$N_5(A)=I_5(A,1)$, $A\wedge B=N_5(I_5(A,N_5(B)))$, and
$\exists_a B=N_5(U_5(a,N_5(B)))$.

The separate-head calculus $\mathcal C_5$ has modus ponens and theorem
generalization $B\mapsto U_5(i,B)$. Its six universally instantiated heads
are
\[
\begin{aligned}
 L_1(D,H)&=(\exists_D U_5(D,H))\to H,\\
 L_2(D,A,H)&=U_5(D,U_5(A,H))\to U_5(A,U_5(D,H)),\\
 L_3(D,A,H)&=M_5(H,D)\to(E_5(H,A)\to M_5(A,D)),\\
 L_4(D,A,H)&=E_5(H,D)\to(M_5(A,H)\to M_5(A,D)),\\
 L_5(D,H)&=\exists_D E_5(H,D),\\
 L_6(D,A,H)&=U_5(D,A)\to(U_5(D,A\to H)\to U_5(D,H)).
\end{aligned}
\]
Here every displayed implication is $I_5$. The remaining heads are
\[
 \rho_5(D,A,H,u)=((H\to u)\to D)\to((D\to H)\to(A\to H)),
 \quad u\le D,
 \qquad \Gamma_5(B)=B\to U_5(B,B),
\]
and the guarded gate $G_5(a,b,c)$ for $c\ne a,b$. Explicitly,
\[
\begin{aligned}
 V(a,b,c)&=M_5(a,b)\to U_5(c,M_5(a,b)),\\
 P_c&=c\to M_5(c,b),\\
 S(a,b,c)&=U_5\bigl(b,(M_5(a,b)\to P_c)\wedge(M_5(c,b)\to c)\bigr),\\
 Q(a,b,c)&=U_5(c,S(a,b,c)\to N_5(P_c)),\\
 G_5(a,b,c)&=N_5(V(a,b,c)\to Q(a,b,c)).
\end{aligned}
\]
A failed guard emits zero. This is derivable, since the dominated
propositional instance with all four parameters zero has code zero.
Every natural number is permitted as a formula code, so the fifth residue
class must receive an interpretation; it is addressed below.

\subsection{Refutations transfer from coefficient four to coefficient five}

The supplied separate-head coefficient-four calibration constructs a
primitive-recursive translation of PA refutations into proofs of code one
in $\mathcal C_4$. Its completeness argument replaces the old
equality-vacuity interface by diagonal freshness, reverse membership
rewriting and a neighborhood-equality witness. These are the supplied
ordinary mathematical derivations, retained in the logical references.
To obtain $\mathcal C_5$ one needs more than a typographical replacement of
four by five, because formula codes are also used as variable labels.

Define a hereditary recoder $F$ by $F(0)=0$ and, for $n>0$ with unique
decoding $n=4P(a,b)+r$, $0\le r<4$, put
\[
                       F(n)=5P(F(a),F(b))+r.
\]
Both decoded components are smaller than $n$. A recursion storing its
previous values therefore defines $F$ primitive recursively, and PA
proves its defining equations. It fixes $0,1,2,3$ and takes positive
arguments to positive values.

Define $J(0)=0$ and, when $m=5P(a,b)+r>0$, put
\[
 J(m)=\begin{cases}4P(J(a),J(b))+r,&r<4,\\0,&r=4.\end{cases}
\]
The same strict component bounds make this primitive recursive. Induction
on $n$ gives $J(F(n))=n$, hence $F$ is injective. Direct substitution,
including the exceptional case $I_4(0,0)=0$, gives
\[
\begin{aligned}
 F(I_4(A,B))&=I_5(F(A),F(B)),\\
 F(M_4(a,b))&=M_5(F(a),F(b)),\\
 F(U_4(a,B))&=U_5(F(a),F(B)),\\
 F(E_4(a,b))&=E_5(F(a),F(b)).
\end{aligned}
\]
The same $F$ is applied to every numerical variable label. Injectivity
preserves binding, repeated labels and the inequalities $c\ne a,b$ in the
gate guard. In particular,
\[
                         F(\Gamma_4(B))=\Gamma_5(F(B)).
\]
This exact identity is why the variable labels must be recoded as well as
the formula constructors.

The recoder is not order-preserving: $40<44$, whereas $F(40)=75>55=F(44)$.
A restricted propositional leaf must therefore be translated using a
proof, not by asserting $F(u)\le F(D)$. The supplied generic proof term
\texttt{M6/Logic/RestrictedImplication.lean:rho8All} derives the
unrestricted schema from MP, the restricted schema and the order rules
\[
 u\le u,\qquad u\le a\Longrightarrow u\le I_5(a,b),\qquad
 u\le b\Longrightarrow u\le I_5(a,b).
\]
These rules hold because $a,b\le P(a,b)\le5P(a,b)$. The proof term is a
fixed finite template, with every restricted leaf justified by those
bounds; it assumes no numerical monotonicity of $F$.

Recursing on a $\mathcal C_4$ proof, recode each ordinary rule and replace
each restricted propositional leaf by this unrestricted template at its
recoded arguments. PA verifies the resulting primitive-recursive proof
transformation. Since $F(1)=1$, composing it with the supplied
PA-to-$\mathcal C_4$ refutation translation proves in PA that every PA
refutation yields a $\mathcal C_5$ proof of one.

\subsection{Soundness in PA, including residue four}

Define a total syntactic translation $T$ into arithmetic formulas. Put
$T(0)=(0=0)$; for $n=5P(a,b)+r>0$ use
\[
\begin{array}{c|l}
 r&T(n)\\\hline
 0&T(a)\to T(b)\\
 1&\operatorname{Bit}(x_a,x_b)\\
 2&\forall x_b\,T(a)\\
 3&x_a=x_b\\
 4&0=0.
\end{array}
\]
The bit predicate expresses that its second argument is a set bit of its
first argument. Its internal auxiliary variables are renamed to prevent
capture. Decoded components are strictly smaller than $n$, so this is a
primitive-recursive translation of syntax. Residue four is interpreted as
closed truth even when such a number occurs as an arbitrary command
parameter or subformula.

Every source variable label appearing in $T(B)$ is below $B$; the truth
and residue-four cases have no such labels. Thus $x_B$ is not free in
$T(B)$, and every diagonal head has a uniformly generated PA proof.
The implication-code identity is a syntactic identity except when both
arguments are zero. In that case the required equivalence is the fixed
propositional theorem
\[
 (0=0)\ \longleftrightarrow\ ((0=0)\to(0=0)).
\]
This handles $I_5(0,0)=0$ without a false disjoint-constructor assumption.

PA proves $\operatorname{Bit}(s,z)\to z<s$ and hence refutes
$\operatorname{Bit}(z,z)$. For every arithmetic formula $\theta$ it also
proves finite binary comprehension:
\[
 \forall n\,\exists s<2^n\,\forall z<n\,
       (\operatorname{Bit}(s,z)\leftrightarrow\theta(z,\vec u)).
\]
A proof is supplied by induction on $n$: retain the old code or add $2^n$
according to excluded middle on $\theta(n,\vec u)$. Inserting the syntax
of $\theta$ into this proof pattern is a primitive-recursive producer of
PA proofs. It is not a decision procedure for $\theta$ or a uniform truth
predicate.

For the guarded gate set $\phi=T(c)$. When $a\ne b$, use the finite
comprehension matrix
\[
 (\operatorname{Bit}(x_a,z)\lor z=x_b)\wedge\phi[z/x_b]
\]
with bound $x_a+x_b+1$. When $a=b$, use instead
$(z=x_b)\wedge\phi[z/x_b]$ with bound $x_b+1$.
In the latter case the first membership antecedent inside the binder is
$\operatorname{Bit}(z,z)$, which PA refutes. In each case the chosen set
satisfies the separation matrix and distinguished-point condition.
Freshness of the label $c$, and capture-avoiding substitution at the
parameter labels, ensure that choosing the set does not alter $\phi$.
This supplies a PA proof of $\neg T(Q(a,b,c))$; the guarded membership
vacuity proves $T(V(a,b,c))$, and propositional reasoning gives $T(G_5)$.
The failed-guard output has its fixed truth proof.

The six remaining heads have ordinary first-order proof templates under
this interpretation; the propositional head is a tautology; MP uses the
implication-code equivalence; theorem generalization uses the corresponding
PA rule. Recursion on a finite $\mathcal C_5$ derivation consequently
produces a PA proof of its translation. Since
$T(1)=\operatorname{Bit}(x_0,x_0)$ has a fixed PA refutation, a
$\mathcal C_5$ proof of one yields a PA refutation. This combines with the
preceding recoding argument to establish, inside PA, equivalence of the two
refutation-existence predicates.

\subsection{The exact controller and its proof compiler}

The command store is $(t,w,d,h)$. The ordinary Cantor stack is
$[a;w]=P(a,w)$, with total projections
$\operatorname{hd}(0)=\operatorname{tl}(0)=0$. In particular $[0;0]=0$.
For $b=d+x$, the changed controller is
\[
\begin{aligned}
 \mathsf S(t,w,d,h)&=
  \begin{cases}\mathsf{halt},&h=1,\\
    (h,\operatorname{tl}w,d,\operatorname{hd}w),&h\ne1,
  \end{cases}\\
 \mathsf{MP}_x(t,w,d,h)&=
  \begin{cases}(b,w,b,0),&I_5(t,b)=h,\\
                 (0,w,b,0),&I_5(t,b)\ne h.
  \end{cases}
\end{aligned}
\]
Successful MP leaves the stack unchanged. This fact is reflected in the
compiler, which reserves its own dummy zero before compiling the premises.

For $D=d+x,H=h+y,A=c$, AX emits the guarded gate on top, clears $d,h$,
and pushes ten entries. In top-to-bottom order the nine generated entries
are
\[
 (q_1,\ldots,q_9)=(\Gamma_5(A),L_6,L_5,L_4,L_3,L_2,L_1,
                          \rho_5(D,A,H,d),U_5(d,h)),
\]
followed by the saved incoming $t$. All generated entries and the gate
have residue zero modulo five except $q_9$, which has residue two and is
at least two. None is one. The complete free-algebra emitter identity,
not just these residues, is part of the delivered certificate.

Writing $A_0=\mathsf{AX}_{0,0,0}$ and executing words left to right, define
\[
 \mathsf{Push}=A_0\mathsf S^9\mathsf{MP}_0,\quad
 \mathsf{Pop}=A_0\mathsf S^{11},\quad
 \mathsf{Discard}=\mathsf{MP}_0,\quad
 \mathsf{Offset}_i=\mathsf{MP}_i.
\]
The ten-entry AX identity gives
\[
\begin{aligned}
 \mathsf{Push}(t,w,d,h)&=(0,[t;w],0,0),\\
 \mathsf{Pop}(t,[a;w],d,h)&=
 \begin{cases}\mathsf{halt},&t=1,\\(t,w,0,a),&t\ne1,\end{cases}\\
 \mathsf{Discard}(t,w,0,h)&=(0,w,0,0),\\
 \mathsf{Offset}_i(0,w,0,0)&=(0,w,i,0).
\end{aligned}
\]
Push compares an implication of residue zero with a universal head of
residue two. Pop tests the saved top on its eleventh Shift. For Discard,
the candidate conclusion is zero, so both MP branches give the same
result, even with arbitrary old cache. Offset includes $i=0$. These are
identities of numerical stacks and remain valid when zero pushes collapse.

For $1\le j\le7$, the selector
\[
 \mathsf{Discard}\,\mathsf{AX}_{A,D,H}\,
                  \mathsf S^{j+1}\mathsf{Pop}^{9-j}
\]
returns $q_j$, preserves the arbitrary underlying stack and clears the
caches. The gate uses
$\mathsf{Discard}\,\mathsf{AX}_{c,b,a}\,\mathsf{Pop}^{10}$.
For an actual restricted propositional leaf with $u\le D$, use
\[
 \mathsf{Discard}\,\mathsf{Offset}_u\,
        \mathsf{AX}_{A,D-u,H}\,\mathsf S^9\mathsf{Pop}.
\]
The subtraction is on target-radix parameters satisfying that leaf's own
valid inequality. No identity relating recoding to subtraction is used.
Generalization is
\[
 \mathsf{Gen}_i=\mathsf{Push}\,\mathsf{Offset}_i\,\mathsf S\,
                       \mathsf{AX}_{i,0,0}\,\mathsf S^{10},
\]
which returns $U_5(i,a)$ from $(a,w,0,0)$, including $a=1$.

The MP suffix $\mathsf S\mathsf S\mathsf{MP}_b\mathsf{Pop}$ sends
$(0,[a,I_5(a,b),0;w],0,0)$ to $(b,w,0,0)$ when $a,b\ne1$ and halts
when $a=1$ or $b=1$. The two Shifts expose the premises; MP leaves
$(b,[0;w],b,0)$; Pop consumes the reserved zero and clears the offset,
or detects a false conclusion. This argument also includes the simultaneous
aliases $I_5(0,0)=0$ and $[0;0]=0$.

Compile axioms by the selectors and generalization by $\mathsf{Gen}$.
At an MP node with minor-proof $p_0$ of $a$ and major-proof $p_1$ of
$I_5(a,b)$, use
\[
 W_p=\mathsf{Discard}\,\mathsf{Push}\,
 W_{p_1}\,\mathsf{Push}\,W_{p_0}\,\mathsf{Push}\,
 \mathsf S\mathsf S\mathsf{MP}_b\mathsf{Pop}.
\]
Induction on the proof tree shows that from every $(t,w,0,0)$ this word
either halts or returns $(\operatorname{concl}p,w,0,0)$.
If $p$ proves one, then $A_0W_p\mathsf{Pop}$ halts from every incoming
store. The leading AX clears arbitrary offsets and cache; the final Pop
exposes and tests the proved one before using the old stack. A plain final
Shift is insufficient, since $\mathsf S(1,0,0,0)=(0,0,0,0)$.

For soundness, maintain derivations in $\mathcal C_5$ for $t,h$ and all
entries of the canonical finite stack. The virtual zero tail receives
its fixed derivation. AX introduces the specified axioms and generalizes
the old cache; successful MP combines the annotated premises; failed MP
returns truth; Shift moves annotations and halts only on an annotated
one. No derivability condition is imposed on the numerical offset $d$.
All annotation updates and proof compilation are primitive recursive,
and their correctness follows by induction in PA.

\subsection{Exhaustive scheduling and the literal-machine endpoint}

The scheduler retains its ordinary Cantor encoding; it does not use the
factor five. Its command encoders are
\[
 B_{\mathsf S}(z)=P(0,z),\quad
 B_{\mathsf{MP}_x}(z)=P(1,P(x,z)),\quad
 B_{\mathsf{AX}_{c,x,y}}(z)=P(c+2,P(y,P(x,z))).
\]
For continuation words, $E([])=0$ and $E(C::V)=1+B_C(E(V))$.
A fresh candidate decodes its value directly; a positive continuation
first subtracts one. Given a universal halting word $W$, put $z=E(W)$.
The fresh candidates
\[
                    n_k=P(k+2,P(0,P(0,z)))
\]
execute $\mathsf{AX}_{k,0,0}$ followed by $W$. They are increasing and
unbounded. The fresh counter begins at two. Every earlier trial either
halts or finishes, since its positive continuation strictly decreases and
all source arithmetic loops have primitive-recursive termination bounds.
Finite induction therefore reaches some selected $n_k$ unless an earlier
halt occurs. The carried store causes no difficulty because $W$ halts
from every store.

The source binding is exact. Relative to the supplied coefficient-four H
source, each of the 61 multiplier blocks changes from four increments to
five; all 3,409 outside instruction sites and all their relocated successors
are preserved. The changed loop has invariant
$T=T_0+5i$, $Y=Y_0-i$, and terminates after $7Y_0+1$ primitive steps
with $T=T_0+5Y_0,Y=0$. It changes formula Cons to
\[
            (t,[a;w],0,0)\longmapsto(I_5(a,t),w,0,0)
\]
in coordinates $(T,W,X,Y)$, while leaving ordinary stack pairing and
scheduler coding unchanged. The complete source has 3,714 sites and
4,815 branch-sensitive outcomes. The no-push MP and cache-only Shift
subtrees are exactly those of the supplied H source, as checked by restoring
the multiplier blocks and comparing the entire resolved graph.

The literal table certificates implement every source outcome by a finite
positive sequence of physical steps, for arbitrary register values,
affine header gap and nonempty spare tail. They preserve the full frame
and untouched tape exterior and admit no premature halt. The complete blank
initialization reaches source entry 25 with logical registers
$(2,0,0,0,0,0,0,0,0)$ after 58,088 physical steps.

A halting source trace is converted to a physical halting trace by
concatenating startup and the certified finite segments. Conversely,
if the physical run first halted at time $T$ but the source had not
halted, simulate the first $T+1$ source steps. Their positive durations
would provide a nonhalting physical prefix longer than $T$, contradicting
determinism. These are finite primitive-recursive trace constructions;
PA proves their correctness by ordinary induction.

The source annotation yields a $\mathcal C_5$ refutation from a source
halt, and arithmetic translation yields a PA refutation. In the other
direction, a PA refutation is recoded into a $\mathcal C_5$ refutation,
compiled into an arbitrary-store halting word, reached by the scheduler,
and physically simulated. This proves (H5), and hence
\[
 \boxed{\mathrm{PA}\vdash
       \neg H_M\ \longleftrightarrow\ \operatorname{Con}(\mathrm{PA}).}
\]
Consistency of PA is required only for the additional external conclusion
that the delivered machine actually does not halt.
